% !TeX program = pdflatex
%% ---------------------------------------------------------------
%%  Korrekturnotiz zu Korollar 2.1.2 (Kapitel A-2, Abschnitt 2.1)
%%  Eigenstaendig uebersetzbar: pdflatex cor-a2-1-2-korrektur.tex
%% ---------------------------------------------------------------
\documentclass[11pt,a4paper]{article}

\usepackage[T1]{fontenc}
\usepackage[utf8]{inputenc}
\usepackage[ngerman]{babel}
\usepackage{lmodern}
\usepackage[margin=3cm]{geometry}
\usepackage{amsmath,amssymb}
\usepackage{mathtools}
\usepackage{csquotes}
\usepackage{enumitem}
\usepackage{verbatim}
\usepackage[colorlinks=true,linkcolor=black,urlcolor=blue]{hyperref}

%% -- Makros aus dem Buchprojekt (hier lokal nachgebildet) --------
\newcommand{\LE}{\mathcal{L}(E)}
\newcommand{\N}{\mathbb{N}}
\newcommand{\R}{\mathbb{R}}
\providecommand{\coloneq}{\mathrel{\vcentcolon=}}
\newcommand{\dr}{\,\mathrm{d}r}
\newcommand{\ds}{\,\mathrm{d}s}

\title{Zum Beweis von \enquote{(d) $\Rightarrow$ (a)}\\
       in Korollar 2.1.2}
\author{Claude \& Ulrich Groh}
\date{23.~Oktober 2026}

\begin{document}
\maketitle

\section*{Der fragliche Satz}

Im Beweis von Korollar~2.1.2 (Kapitel A-2, Abschnitt~2.1) steht am
Beginn des letzten Beweisteils:

\begin{quote}
It remains to show that (d) implies (a).
Assertion (d) and the semigroup property imply that for all $t \geq 0$
there exist bounded operators $T(t) \in \LE$ such that
\[
    u(t,f) = T(t)f
    \quad \text{if} \quad
    f \in D(A) .
\]
Moreover,
\[
    \sup_{0 \leq t \leq 1} \|T(t)\| < \infty .
\]
\end{quote}

\noindent
Die Formulierung scheint mir an zwei Stellen nicht zu stimmen.

\section*{Was (d) tatsächlich liefert}

Bedingung (d) lautet: $A$ erfüllt (EU), hat dichten Definitionsbereich,
und für jede Folge $(f_n)$ in $D(A)$ mit $\lim_{n\to\infty} f_n = 0$ gilt
$\lim_{n\to\infty} u(t,f_n) = 0$ \emph{gleichmäßig in} $t \in [0,1]$.

Daraus folgt der Reihe nach:

\begin{enumerate}[label=(\roman*)]
\item Aus der Eindeutigkeit in (EU) ergibt sich, dass
      $f \mapsto u(\cdot,f)$ auf $D(A)$ linear ist.
\item Die lineare Abbildung
      $\eta \colon D(A) \to C([0,1],E)$, $\eta(f) = u(\cdot,f)$,
      ist nach (d) folgenstetig im Nullpunkt bezüglich der Norm von $E$,
      also beschränkt. Es gibt somit ein $M \geq 0$ mit
      \[
          \sup_{0 \leq t \leq 1} \|u(t,f)\| \leq M\|f\|
          \qquad (f \in D(A)).
      \]
\item Da $D(A)$ dicht in $E$ ist, setzt sich $f \mapsto u(t,f)$ für jedes
      $t \in [0,1]$ eindeutig zu einem Operator $T(t) \in \LE$ fort, und
      zugleich erhält man $\sup_{0 \leq t \leq 1}\|T(t)\| \leq M < \infty$.
\end{enumerate}

\section*{Zwei Einwände}

\paragraph{1.\ Der Bereich $t \geq 0$ ist zu groß.}
Die obige Argumentation liefert die Operatoren $T(t)$ nur für
$t \in [0,1]$; für $t > 1$ gibt (d) nichts her. Genau deshalb wird
$T(t)$ im übernächsten Absatz des Beweises ja erst durch
$T(t) \coloneq T(1)^{n}\,T(s)$ mit $t = n+s$, $s \in [0,1)$,
\emph{definiert}. Mit der bisherigen Formulierung ist dieser Absatz
redundant bzw.\ widersprüchlich: $T(t)$ wäre für $t > 1$ bereits
vorhanden und würde anschließend nochmals festgelegt.

\paragraph{2.\ \enquote{the semigroup property} steht noch nicht zur Verfügung.}
Eine Halbgruppe gibt es an dieser Stelle noch gar nicht --- sie wird
erst konstruiert. Benötigt werden die Eindeutigkeit der Lösungen (für
die Linearität) und die Stetigkeitsbedingung aus (d) (für die
Beschränktheit). Der Verweis auf die Halbgruppeneigenschaft ist daher
auch sachlich schief.

Außerdem sollte am Schluss des Absatzes die Dichtheit von $D(A)$
erwähnt werden, damit die Halbgruppeneigenschaft von $D(A)$ auf ganz
$E$ übergeht.

\section*{Vorschlag für die Neufassung}

\begin{quote}
It remains to show that (d) implies (a).
Assertion (d) and the uniqueness of the solutions imply that for every
$t \in [0,1]$ there exists a bounded operator $T(t) \in \LE$ such that
\[
    u(t,f) = T(t)f
    \quad \text{if} \quad
    f \in D(A),
\]
and that
\[
    \sup_{0 \leq t \leq 1} \|T(t)\| < \infty .
\]
In fact, by the uniqueness of the solutions the mapping
$f \mapsto u(\cdot,f)$ from $D(A)$ into $C([0,1],E)$ is linear, and by
(d) it is continuous at $0$, hence bounded. Since $D(A)$ is dense in
$E$, it extends uniquely to a bounded operator from $E$ into
$C([0,1],E)$.
It follows that $T(\cdot)f$ is strongly continuous on $[0,1]$ for all
$f \in E$ (since it is so for $f \in D(A)$ and $D(A)$ is dense).

Let $t > 1$.
There exist unique $n \in \N$ and $s \in [0,1)$ such that $t = n+s$.
Let
\[
    T(t) \coloneq T(1)^{n} \cdot T(s) .
\]
From the uniqueness of the solutions it follows that $T(t)f = u(t,f)$
for all $t \geq 0$ and $f \in D(A)$, as well as
$T(t+s)f = T(s)T(t)f$ for all $f \in D(A)$ and $s$, $t \geq 0$.
Since $D(A)$ is dense, $T$ is a semigroup.
\end{quote}

\section*{Bemerkung zum Gruppenfall}

Im Beweis von Theorem~2.1.6 steht dieselbe Konstruktion, dort aber
korrekt in der Form \enquote{If (b) holds then there exist operators
$T(t) \in \mathcal{L}(E)$ \ldots}. Dort wird zusätzlich der Satz von
Hille--Phillips über die Messbarkeit benötigt, weil die Konvergenz in
der Voraussetzung nur punktweise in $t$ gefordert wird. Im
Halbgruppenfall ist die Gleichmäßigkeit auf $[0,1]$ Teil der
Voraussetzung (d); deshalb geht es hier direkter.

\appendix
\section*{Anhang: Quelltext des Vorschlags}
\addcontentsline{toc}{section}{Anhang}

\small
\begin{verbatim}
It remains to show that (d) implies (a). 
Assertion~(d) and the uniqueness of the solutions imply that for 
every $ t \in [0,1] $ there exists a bounded operator $ T(t) \in \LE $ 
such that 
%
\[
	u(t,f) = T(t)f 
	\quad \text{if} \quad 
	f \in D(A) ,
\]
%
and that 
%
\[
	\sup_{0 \leq t \leq 1} \|T(t)\| < \infty .
\]
%
In fact, by the uniqueness of the solutions the mapping 
$ f \mapsto u(\cdot,f) $ from $ D(A) $ into $ C([0,1],E) $ is linear, 
and by~(d) it is continuous at $ 0 $, hence bounded. 
Since $ D(A) $ is dense in $ E $, it extends uniquely to a bounded 
operator from $ E $ into $ C([0,1],E) $.
It follows that $ T(\cdot)f $ is strongly continuous on $ [0,1] $ for 
all $ f \in E $ (since it is so for $ f \in D(A) $ and $ D(A) $ is dense).

Let $ t > 1 $.
There exist unique $ n \in \N $ and $ s \in [0,1) $ such that $ t = n + s $.
Let 
%
\[
	T(t) \coloneq T(1)^{n} \cdot T(s) .
\]
%
From the uniqueness of the solutions it follows that $ T(t)f = u(t,f) $ 
for all $ t \geq 0 $ and $ f \in D(A) $, as well as 
$ T(t+s)f = T(s)T(t)f $ for all $ f \in D(A) $ and $ s$, $t \geq 0 $.
Since $ D(A) $ is dense, $ T $ is a semigroup.
\end{verbatim}

\end{document}
